% %--------------------------
% %	CONFIGURATION
% %--------------------------
\documentclass[11pt,a4paper]{moderncv}
\moderncvstyle{classic}
\moderncvcolor{black}
\definecolor{firstnamecolor}{rgb}{0,0,0}

% Basic packages
\usepackage[utf8]{inputenc}
\usepackage[T1]{fontenc}
\usepackage[scale=0.85]{geometry}
\usepackage[dvipsnames]{xcolor}
\definecolor{DarkBlue}{RGB}{0, 0, 150}
\usepackage[colorlinks=true, urlcolor=DarkBlue, linkcolor=DarkBlue]{hyperref}

\usepackage{microtype}

% Configure font settings
\renewcommand{\rmdefault}{ppl}
\renewcommand{\sfdefault}{phv}
\renewcommand*{\namefont}{\fontsize{26}{18}\mdseries}

% Configure microtype
\microtypesetup{expansion=false}

% Load sensitive information
\input{../../other_material/personal_info.tex}

% Title document
\title{Curriculum Vitae}

% %--------------------------

\begin{document}
\makecvtitle

\section*{Introduction}
I am an <<ADJECTIVE>> Data Scientist with <<YEARS_OF_EXPERIENCE>> years of experience applying ML across <<RELEVANT_INDUSTRIES>>. I have a strong mathematical background, a perfect GPA from my Master's in Statistical Learning and Data Analytics, and hands-on experience with <<RELEVANT_TECHNICAL_THEMES>>. After some time in industry, I am excited to return to academia to work on more fundamental research in <<TARGET_RESEARCH_AREA>>.

\section{Education}

\cventry{2015--2020}{Royal Institute of Technology, KTH -- Data Science and Engineering}{}{GPA: 5.0 (out of 5.0)}{}{Master: Statistical Learning and Data Analytics, Bachelor: Engineering Physics. \href{https://kth.diva-portal.org/smash/get/diva2:1431327/FULLTEXT02.pdf}{See thesis}. \newline Relevant coursework: <<RELEVANT_COURSEWORK_LIST>>.}

\cventry{2017--2022}{Stockholm University, SU - Economics}{Supplementary Bachelor's Degree}{}{}{}

\cventry{2019--2019}{Swiss Federal Institute of Technology, ETH - Applied Mathematics}{Exchange}{}{}{}

% ---------------------------
% Selected Work Experience
% ---------------------------
\section{Selected Work Experience}

% ---------------------------
\large{Data Scientist \& Project Manager @ Signum Life Science}
\normalsize
\begin{itemize}
\item Own, extend, and maintain next-best-action engine for pharma sales, integrating explainable AI.
\item Improve and maintain Denmark-specific price forecasting model for blind pharmaceutical auctions.
\item Organised a company-wide hackathon with 40 participants, fostering innovation across teams.
\end{itemize}

{\footnotesize\textit{Keywords: Explainable AI, AWS, Azure Databricks, forecasting, XGBoost, pricing, ML engineering, project management}}

\hfill{\small{\textit{\hyperref[sec:signum]{Read more...}}}}

% ---------------------------
\large{Principal Data Scientist @ Valcon}
\normalsize
\begin{itemize}
    \item Led design and development of a real-time multi-modal ML pipeline deployed using Azure ML.
    \item Mentored junior colleagues and held trainings in xAI and Git.
    \item Won "People's Choice Award" for exemplary work, leadership, team spirit, and curiosity.
\end{itemize}

{\footnotesize\textit{Keywords: Real-time ML, NLP, explainability in AI, MLOps, multi-modal, Azure, Docker, mentoring, stakeholder management}}

\hfill{\small{\textit{\hyperref[sec:valcon]{Read more...}}}}

% ---------------------------
\large{Researcher within Data Science @ RaySearch Laboratories}
\normalsize
\begin{itemize}
    \item Designed a 3D computer-vision pipeline for automated cancer treatment.
    \item Built lexicographic dose-optimisation for palliative care. Continuously validated with clinicians.
\end{itemize}

{\footnotesize\textit{Keywords: 3D computer vision, probabilistic modelling, latent representation, segmentation, autoencoders, optimisation, healthcare AI}}

\hfill{\small{\textit{\hyperref[sec:raysearch]{Read more...}}}}

% ---------------------------
% Solo Side Project
% ---------------------------
\section{Solo Side Project}
\normalsize
Developing an end-to-end CV system to improve e-commerce navigation by making product images interactive. Works by creating embeddings of products to make them searchable across the webshop.
\begin{itemize}
  \item Web scraping pipeline to dynamically index all items on a webshop.
  \item Object detection and classification using YOLOv8 and OpenClip to extract all items.
  \item Chromium plugin to display interactive links, improving user navigation.
\end{itemize}
{\footnotesize\textit{Keywords: Computer vision, object detection, segmentation, semantic embedding, full-stack ML engineering, fine-tuning, web scraping}}
\hfill{\small{\textit{\hyperref[sec:iris]{Read more…}}}}

% ---------------------------
% Skills and Other
% ---------------------------
\section{Skills and Other}
\normalsize
\cvitem{Technology}{Python (PyTorch, TensorFlow, SHAP, YOLO), C++, Matlab, MS Azure (Data Science Associate Certified), Databricks, AWS, MongoDB, Qdrant, Docker, 3D Printing, Fusion 360.}
\cvitem{Language}{\textit{Native}: English, Swedish. \textit{Intermediate}: Danish.}
\cvitem{Personal}{In my spare time I enjoy working on projects, such as designing and building furniture.}
\cvitem{Other}{Member of Mensa Sweden.}

\newpage

% ---------------------------
% Detailed Work Experience
% ---------------------------
\section{Detailed Work Experience}

% ---------------------------
\phantomsection\label{sec:signum}
\cventry{2024--Present}{Data Scientist and Project Manager}{Signum Life Science}{Copenhagen}{}{
    At Signum, I maintain and lead development on Kwarts, a next-best-action engine for pharmaceutical sales. Originally built externally, I now own the system, extending it with features like explainable AI in collaboration with UX designers, and supporting its commercialisation by aiding our sales representatives. Kwarts adapts to varied customer data maturity, from rule-based logic to deep learning, and is deployed in AWS for scalability.
    \newline\hspace*{2em}
    I also maintain and improve our pharmaceutical price forecasting model, which helps clients navigate Denmark's fortnightly blind auctions. After stabilising the XGBoost-based system, I am now driving efforts toward a more robust and modern implementation.
    \newline\hspace*{2em}
    Beyond model development, I have also organised an internal week-long hackathon that included up to 40 participants. The event fostered innovation and collaboration, resulting in several new ideas and prototypes.
}{}

% ---------------------------
\phantomsection\label{sec:valcon}
\cventry{2022--2024}{Principal Data Scientist}{Valcon}{Copenhagen}{}{
    Led data science projects across pharma, energy, IoT, and manufacturing. Mentored junior colleagues, led xAI and Git trainings, and shaped internal practices. Regularly presented to non-technical stakeholders, honing my ability to translate complex ideas clearly.
    \newline\hspace*{2em}
    One project involved training a real-time deep neural network with an advanced custom training pipeline and modular architecture. The training data consisted of natural-language and high dimensional multi-class labels which made for a challenging prediction problem. I was the lead data scientist and architect of the solution and responsible for implementing the end-to-end in-production pipeline.
    \newline\hspace*{2em}
    I was awarded the "Valcon People's Choice Award" for exemplary work, leadership, team spirit, and scientific curiosity.
}{}

% ---------------------------
\phantomsection\label{sec:valtech}
\cventry{2021--2022}{Data Scientist}{Valtech}{Copenhagen}{}{
    Delivered ML solutions in e-commerce. Built time series forecasts using Facebook Prophet, and applied clustering to behavioural data to support personalisation. Worked closely with business stakeholders to ensure actionable outputs.
}{}

% ---------------------------
\phantomsection\label{sec:raysearch}
\cventry{2019--2021}{Researcher within Data Science}{RaySearch Laboratories}{Stockholm}{}{
    At RaySearch Laboratories, I worked on automating radiotherapy planning using a combination of deep learning, probabilistic modelling, and multi-objective optimisation. I developed a pipeline that extracted latent anatomical features from segmented 3D CT scans using autoencoders, and used these to identify similar patients via a learned similarity metric. Dose-volume histograms and spatial dose distributions were then predicted using sparse Gaussian processes, providing a probabilistic framework for dose planning conditioned on patient geometry. This enabled personalised, data-driven dose estimation and laid the groundwork for more robust automation in treatment.
    \newline\hspace*{2em}
    Alongside this, I redesigned some of RaySearch's lexicographic optimisation algorithms to better reflect clinical decision hierarchies, particularly in palliative care where trade-offs between tumour control and organ sparing are nuanced. I collaborated closely with medical physicists and clinicians throughout, ensuring the models aligned with both oncological standards and practical workflow requirements.
}{}

% ---------------------------
% Other Work Experience
% ---------------------------
\section{Other Detailed Experience}

% ---------------------------
\phantomsection\label{sec:iris}
\cventry{2025-Present}{Solo Project}{Project iris}{}{}{
    In this personal project I am building an end-to-end computer vision system for making product images on e-commerce sites interactive. The pipeline detects individual items using a YOLOv8 model, and embeds them semantically using OpenCLIP. The embeddings are then compared to support intelligent product linking across all images and products in a web store.
    \newline\hspace*{2em}
    The system includes a full-stack setup with a FastAPI backend, MongoDB for persistence, a Qdrant index for vector search, and a custom-built Chromium plug-in that overlays clickable links directly onto web shop imagery in real-time. The project also includes a dynamic scraping and indexing pipeline for ingesting and structuring product data from arbitrary e-commerce sites.
    \newline\hspace*{2em}
    The project has been a learning opportunity for fine-tuning of deep learning models and practical ML engineering.
}{}

\end{document}